% GENERATED FILE. Edit canonical lessons, then run npm run generate:books.
% canonical-source-hash: fnv1a64:ed8096df6f96fb8d
% canonical-lessons: ES-C600-ciudadania, ES-C600-colectivo, ES-C600-civismo, ES-C600-ciudadano, ES-C600-pertenecer

\chapter{La sociedad}
\label{ch:es-la-sociedad}
\hlchaptermodality{600}

\begin{chapteropening}
\textbf{By the end of this chapter:} \emph{I can talk about citizenship, civic behaviour and the groups people belong to.}

Complete the last lesson of chapter 600: I can talk about citizenship, civic behaviour and the groups people belong to.
\end{chapteropening}

\section[la ciudadanía]{la ciudadanía — the citizenship; the citizens}

\label{lesson:ES-C600-ciudadania}



\begin{quote}
Before the new one: say the Spanish for the conclusion, then the Spanish for the method.
\end{quote}

\subsection*{You'll want to know: la ciudadanía}
\textbf{la ciudadanía} — \textquotedblleft{}the citizenship; the citizens\textquotedblright{}.

\begin{quote}
Tiene la ciudadanía española.
\end{quote}

\begin{quote}
\emph{He has Spanish citizenship.}
\end{quote}

It also means all the citizens together: \emph{la ciudadanía decide}.

\begin{cousinweb}
\emph{ciudadanía} is built on \emph{ciudadano}, from \emph{ciudad}, inherited from Latin \emph{civitatem} 'city, state' (EtymDB).
\end{cousinweb}

\begin{grammarlens}[title={What you've built}]
One more word for telling, explaining and describing at B1.
\end{grammarlens}

\subsection*{Your turn}
\begin{itemize}
\raggedright
  \item \emph{Say it:} \emph{la ciudadanía}
  \item \emph{Say it:} \emph{la ciudadanía}, once more
  \item \emph{Say it:} a sentence of your own with \emph{la ciudadanía}
  \item \emph{Recall:} say the Spanish for the conclusion, then the Spanish for the method, then say \emph{la ciudadanía} again
\end{itemize}

\begin{tcolorbox}[breakable,colback=teal!4,colframe=teal!35!black,title={Before you move on}]
What does la ciudadanía mean? (\textquotedblleft{}The citizenship; the citizens\textquotedblright{}.) Say it once more.
\end{tcolorbox}
\section[el colectivo]{el colectivo — the group; collective}

\label{lesson:ES-C600-colectivo}



\begin{quote}
Before the new one: say the Spanish for the method, then the Spanish for the citizenship.
\end{quote}

\subsection*{You'll want to know: el colectivo}
\textbf{el colectivo} — \textquotedblleft{}the group; collective\textquotedblright{}.

\begin{quote}
Es un colectivo de artistas.
\end{quote}

\begin{quote}
\emph{It is a collective of artists.}
\end{quote}

Also a group of people with something in common: \emph{el colectivo de jóvenes}. In Argentina \emph{el colectivo} is also the city bus.

\begin{cousinweb}
\emph{colectivo} is a learned borrowing from Latin \emph{collectivus} 'gathered together', built on \emph{colligere}, from \emph{legere} 'to gather' (Lewis \& Short).
\end{cousinweb}

\begin{grammarlens}[title={What you've built}]
One more word for telling, explaining and describing at B1.
\end{grammarlens}

\subsection*{Your turn}
\begin{itemize}
\raggedright
  \item \emph{Say it:} \emph{el colectivo}
  \item \emph{Say it:} \emph{el colectivo}, once more
  \item \emph{Say it:} a sentence of your own with \emph{el colectivo}
  \item \emph{Recall:} say the Spanish for the method, then the Spanish for the citizenship, then say \emph{el colectivo} again
\end{itemize}

\begin{tcolorbox}[breakable,colback=teal!4,colframe=teal!35!black,title={Before you move on}]
What does el colectivo mean? (\textquotedblleft{}The group; collective\textquotedblright{}.) Say it once more.
\end{tcolorbox}
\section[el civismo]{el civismo — civic-mindedness, good citizenship}

\label{lesson:ES-C600-civismo}



\begin{quote}
Before the new one: say the Spanish for the citizenship, then the Spanish for the group.
\end{quote}

\subsection*{You'll want to know: el civismo}
\textbf{el civismo} — \textquotedblleft{}civic-mindedness, good citizenship\textquotedblright{}.

\begin{quote}
Tirar basura en la calle es falta de civismo.
\end{quote}

\begin{quote}
\emph{Dropping litter in the street shows a lack of civic-mindedness.}
\end{quote}

Respect for shared spaces and for the people around you.

\begin{cousinweb}
\emph{civismo} is built on \emph{cívico}, a learned borrowing from Latin \emph{civicus} 'of citizens' (EtymDB), from \emph{civis} 'citizen'; it probably came in through French \emph{civisme}.
\end{cousinweb}

\begin{grammarlens}[title={What you've built}]
One more word for telling, explaining and describing at B1.
\end{grammarlens}

\subsection*{Your turn}
\begin{itemize}
\raggedright
  \item \emph{Say it:} \emph{el civismo}
  \item \emph{Say it:} \emph{el civismo}, once more
  \item \emph{Say it:} a sentence of your own with \emph{el civismo}
  \item \emph{Recall:} say the Spanish for the citizenship, then the Spanish for the group, then say \emph{el civismo} again
\end{itemize}

\begin{tcolorbox}[breakable,colback=teal!4,colframe=teal!35!black,title={Before you move on}]
What does el civismo mean? (\textquotedblleft{}Civic-mindedness, good citizenship\textquotedblright{}.) Say it once more.
\end{tcolorbox}
\section[el ciudadano]{el ciudadano — the citizen}

\label{lesson:ES-C600-ciudadano}



\begin{quote}
Before the new one: say the Spanish for the group, then the Spanish for civic-mindedness.
\end{quote}

\subsection*{You'll want to know: el ciudadano}
\textbf{el ciudadano} — \textquotedblleft{}the citizen\textquotedblright{}.

\begin{quote}
Todos los ciudadanos pueden votar.
\end{quote}

\begin{quote}
\emph{All citizens can vote.}
\end{quote}

A woman is \emph{la ciudadana}.

\begin{cousinweb}
\emph{ciudadano} is built on \emph{ciudad}, inherited from Latin \emph{civitatem} 'city, state', built on \emph{civis} 'citizen' (EtymDB).
\end{cousinweb}

\begin{grammarlens}[title={What you've built}]
One more word for telling, explaining and describing at B1.
\end{grammarlens}

\subsection*{Your turn}
\begin{itemize}
\raggedright
  \item \emph{Say it:} \emph{el ciudadano}
  \item \emph{Say it:} \emph{el ciudadano}, once more
  \item \emph{Say it:} a sentence of your own with \emph{el ciudadano}
  \item \emph{Recall:} say the Spanish for the group, then the Spanish for civic-mindedness, then say \emph{el ciudadano} again
\end{itemize}

\begin{tcolorbox}[breakable,colback=teal!4,colframe=teal!35!black,title={Before you move on}]
What does el ciudadano mean? (\textquotedblleft{}The citizen\textquotedblright{}.) Say it once more.
\end{tcolorbox}
\section[pertenecer]{pertenecer — to belong}

\label{lesson:ES-C600-pertenecer}



\begin{quote}
Before the new one: say the Spanish for civic-mindedness, then the Spanish for the citizen.
\end{quote}

\subsection*{You'll want to know: pertenecer}
\textbf{pertenecer} — \textquotedblleft{}to belong\textquotedblright{}.

\begin{quote}
Este libro pertenece a la biblioteca.
\end{quote}

\begin{quote}
\emph{This book belongs to the library.}
\end{quote}

Like \emph{conocer}: \emph{pertenezco}. Always with \emph{a}.

\begin{cousinweb}
\emph{pertenecer} is inherited from Latin \emph{pertinere} 'to stretch to, to belong', built from \emph{per} and \emph{tenere} 'to hold', with the ending \emph{-ecer} (EtymDB).
\end{cousinweb}

\begin{grammarlens}[title={What you've built}]
One more word for telling, explaining and describing at B1.
\end{grammarlens}

\subsection*{Your turn}
\begin{itemize}
\raggedright
  \item \emph{Say it:} \emph{pertenecer}
  \item \emph{Say it:} \emph{pertenecer}, once more
  \item \emph{Say it:} a sentence of your own with \emph{pertenecer}
  \item \emph{Recall:} say the Spanish for civic-mindedness, then the Spanish for the citizen, then say \emph{pertenecer} again
\end{itemize}

\begin{tcolorbox}[breakable,colback=teal!4,colframe=teal!35!black,title={Before you move on}]
What does pertenecer mean? (\textquotedblleft{}To belong\textquotedblright{}.) Say it once more.
\end{tcolorbox}
